% Zone „Das kennst du schon“ – aus bank/flaechen/zone.jsonl (zusammenbau v0.1)
\einheitenkopf[zone][Kennst du schon]{Das kennst du schon}
\setcounter{aufgabe}{0}

%% TODO Titel: Ich-kann-Satz fehlt in der Bank (Platzhalter: Kettenname) – Nr. 1
%% TODO Anweisung: ein Satz über der ersten Teilaufgabe fehlt in der Bank – Nr. 1
\begin{aufgabe}{Längeneinheiten umrechnen (mm, cm, m), gemischte Angaben angleichen}
\begin{teile}
\teil $3$ m – wie viele cm? \leerfeld[cm]
\teil $1{,}35$ m – wie viele cm? \leerfeld[cm]
\end{teile}
\end{aufgabe}

%% TODO Titel: Ich-kann-Satz fehlt in der Bank (Platzhalter: Kettenname) – Nr. 2
%% TODO Anweisung: ein Satz über der ersten Teilaufgabe fehlt in der Bank – Nr. 2
\begin{aufgabe}{Flächeneinheiten (cm², m²) und Umrechnungszahl 100}
\begin{teile}
\teil $4$ cm² – wie viele mm²? \leerfeld mm²
\teil $250$ cm² – wie viele dm²? \leerfeld dm²
\end{teile}
\end{aufgabe}

%% TODO Titel: Ich-kann-Satz fehlt in der Bank (Platzhalter: Kettenname) – Nr. 3
%% TODO Anweisung: ein Satz über der ersten Teilaufgabe fehlt in der Bank – Nr. 3
\begin{aufgabe}{Multiplizieren und Dividieren mit Dezimalzahlen (Kommazahl mal Kommazahl, zweistellig geteilt durch einstellig)}
\begin{teile}
\teil $0{,}5 \cdot 8$ – wie viel? \leerfeld
\teil $2{,}5 \cdot 1{,}4$ – wie viel? \leerfeld
\end{teile}
\end{aufgabe}

%% TODO Titel: Ich-kann-Satz fehlt in der Bank (Platzhalter: Kettenname) – Nr. 4
%% TODO Anweisung: ein Satz über der ersten Teilaufgabe fehlt in der Bank – Nr. 4
\begin{aufgabe}{Formel nach einer Größe umstellen (A = a · b → b = A : a)}
\begin{teile}
\teil $P = q \cdot r$ – nach $r$ umstellen?
\teil $G = m \cdot p$ mit $G = 7{,}2$ und $m = 3$ – $p$? p = \leerfeld
\end{teile}
\end{aufgabe}

%% TODO Titel: Ich-kann-Satz fehlt in der Bank (Platzhalter: Kettenname) – Nr. 5
%% TODO Anweisung: ein Satz über der ersten Teilaufgabe fehlt in der Bank – Nr. 5
\begin{aufgabe}{Wurzel ziehen bei Quadratzahlen (√49)}
\begin{teile}
\teil $\sqrt{25}$ – wie viel? \leerfeld
\teil $\sqrt{144}$ – wie viel? \leerfeld
\end{teile}
\end{aufgabe}

%% TODO Titel: Ich-kann-Satz fehlt in der Bank (Platzhalter: Kettenname) – Nr. 6
%% TODO Anweisung: ein Satz über der ersten Teilaufgabe fehlt in der Bank – Nr. 6
\begin{aufgabe}{Rechten Winkel erkennen und einzeichnen (Geodreieck)}
\begin{teile}
\teil Ist das ein rechter Winkel? \janein

\winkel{90}{\alpha}
\teil An welcher Ecke ist ein rechter Winkel? Prüfe mit dem Geodreieck und markiere ihn.

\viereck{(0,0)}{(5,0)}{(5,3)}{(1,2.5)}
\end{teile}
\end{aufgabe}

%% TODO Titel: Ich-kann-Satz fehlt in der Bank (Platzhalter: Kettenname) – Nr. 7
%% TODO Anweisung: ein Satz über der ersten Teilaufgabe fehlt in der Bank – Nr. 7
\begin{aufgabe}{Kreisfläche (π · r²)}
\begin{teile}
\teil Kreis mit Radius $r = 1$ m – Fläche? Runde auf zwei Stellen. \leerfeld[m²]
\teil Kreis mit Radius $r = 1{,}5$ m – Fläche? Runde auf zwei Stellen. \leerfeld[m²]
\end{teile}
\end{aufgabe}

%% TODO Titel: Ich-kann-Satz fehlt in der Bank (Platzhalter: Kettenname) – Nr. 8
%% TODO Anweisung: ein Satz über der ersten Teilaufgabe fehlt in der Bank – Nr. 8
\begin{aufgabe}{Flächeneinheiten (cm², m²) und Umrechnungszahl 100 – Fehler finden}
\begin{teile}
\teil Mia rechnet: $8$ cm² $= 80$ mm². Wo ist der Fehler?
\end{teile}
\end{aufgabe}

%% TODO Titel: Ich-kann-Satz fehlt in der Bank (Platzhalter: Kettenname) – Nr. 9
%% TODO Anweisung: ein Satz über der ersten Teilaufgabe fehlt in der Bank – Nr. 9
\begin{aufgabe}{Flächeneinheiten (cm², m²) und Umrechnungszahl 100 – selbst rechnen}
\begin{teile}
\teil $9$ dm² – wie viele cm²? \leerfeld[cm²]
\end{teile}
\end{aufgabe}
